\chapter{Dualidad, Hahn--Banach y geometría convexa}\label{cap:nuevo10}\label{chap:21}\label{chap:24}

Una funcional lineal asigna un escalar a cada vector y puede interpretarse como una medición: su núcleo registra los vectores que esa medición no detecta. El dual reúne todas las mediciones continuas. Estudiaremos cómo describen hiperplanos y recuperan la norma de un vector, y cómo las cotas inferiores permiten reconocer isomorfismos sobre la imagen.

El puente hacia la geometría será el teorema de Hahn--Banach: extender una funcional conservando una cota permite separar conjuntos convexos. Sus puntos de máximo identificarán después caras extremas, y el teorema de Krein--Milman reconstruirá un compacto convexo a partir de sus puntos extremos mediante combinaciones convexas y límites. Las demostraciones se presentan aquí en el orden de sus dependencias.

Se requieren la norma de operador y su completitud cuando el codominio es completo, las propiedades de los compactos métricos y el lema de Zorn. Este último afirma que un conjunto parcialmente ordenado en el que toda cadena tiene una cota superior posee un elemento maximal. Lo aplicaremos primero a extensiones compatibles de funcionales y después a subconjuntos extremos compactos. En cada aplicación comprobaremos expresamente la condición sobre las cadenas.

\section{Funcionales e hiperplanos}\label{lectura:21}

\begin{definicion}
El \emph{dual continuo} de $V$ es $V^*=\Bcal(V,\F)$. Sus elementos son funcionales lineales continuas y llevan la norma de operador
\[\norm\varphi=\sup_{\norm v\leq1}|\varphi(v)|.\]
El dual algebraico contiene todas las funcionales lineales, sin exigir continuidad, y puede ser mayor que $V^*$.
\end{definicion}
\index{dual!continuo}\index{funcional lineal}

\begin{proposicion}
El espacio $V^*$ es de Banach, aunque $V$ no sea completo.
\end{proposicion}
\begin{proof}
Los cuerpos $\R$ y $\C$ son completos. Aplicamos el teorema~\ref{ch20:banach-operadores} con $W=\F$.
\end{proof}

\begin{proposicion}\label{ch21:kernel-proporcional}
Dos funcionales lineales no nulas tienen el mismo núcleo si y solo si son múltiplos escalares no nulos una de la otra.
\end{proposicion}
\begin{proof}
Si $\psi=c\varphi$ y $c\ne0$, sus conjuntos de ceros coinciden. Recíprocamente, si sus núcleos coinciden, elija $v_0$ con $\varphi(v_0)=1$. Esto es posible: primero elija un vector donde $\varphi$ no se anule y divídalo entre su valor. Para cualquier $v$, el vector $v-\varphi(v)v_0$ pertenece a $\ker\varphi$, pues
\[
 \varphi(v-\varphi(v)v_0)=\varphi(v)-\varphi(v)\cdot1=0.
\]
También pertenece a $\ker\psi$, de modo que $\psi(v)=\psi(v_0)\varphi(v)$. El escalar $\psi(v_0)$ no es cero, porque $v_0$ no está en el núcleo común. Si se admiten funcionales nulas, la igualdad de núcleos igual a $V$ significa que ambas son cero.
\end{proof}

\begin{teorema}\label{ch21:hiperplanos}
Los subespacios propios maximales de un espacio vectorial $V$ son exactamente los núcleos de funcionales lineales no nulas. Cada uno tiene codimensión $1$. La correspondencia es biyectiva si se identifican funcionales que difieren por un factor escalar no nulo.
\end{teorema}
\index{hiperplano}
\begin{proof}
Si $M=\ker\varphi$ y $\varphi(v_0)=1$, todo $v$ admite la descomposición
\[
 v=\underbrace{v-\varphi(v)v_0}_{\in M}+\varphi(v)v_0.
\]
Es única, porque $M\cap\F v_0=\{0\}$. Así $V=M\oplus\F v_0$. Todo subespacio que contiene estrictamente a $M$ contiene algún $w$ con $\varphi(w)\ne0$; la misma descomposición usando $w/\varphi(w)$ muestra que ese subespacio es $V$. Por ello $M$ es maximal.

Recíprocamente, si $M$ es propio maximal, elija $v_0\notin M$. El subespacio $M+\F v_0$ contiene estrictamente a $M$, por lo que es $V$. Como la suma es directa, la regla $\varphi(m+\alpha v_0)=\alpha$ está bien definida, es lineal, no nula y tiene núcleo $M$. La proposición~\ref{ch21:kernel-proporcional} prueba que la única indeterminación es un factor no nulo.
\end{proof}

El criterio de continuidad por cerradura del núcleo ya se demostró en la proposición~\ref{ch21:kernel-cerrado}, al estudiar operadores de rango finito. Aplicado a las funcionales de esta sección, enlaza su descripción algebraica con la topología.

En consecuencia, los hiperplanos cerrados corresponden a las clases de funcionales continuas no nulas. Además, el núcleo de una funcional discontinua es denso: su cierre es un subespacio que contiene al núcleo; por maximalidad, ese cierre es el propio núcleo o todo $V$. La primera posibilidad lo haría cerrado y forzaría continuidad.

\section{Isomorfismos y cotas inferiores}

\begin{definicion}
Un \emph{isomorfismo de espacios normados} $T:V\to W$ es una transformación lineal biyectiva para la cual $T$ y $T^{-1}$ son continuas. Una aplicación lineal es una \emph{isometría lineal} si $\norm{Tv}=\norm v$ para todo $v$; no se exige sobreyectividad. Un isomorfismo isométrico es, además, biyectivo. Una aplicación puede ser un isomorfismo de $V$ sobre su imagen sin ser sobreyectiva sobre $W$.
\end{definicion}
\index{isomorfismo!de espacios normados}

\begin{teorema}\label{ch21:cota-inferior}
Un operador lineal $T:V\to W$ es un isomorfismo de $V$ sobre $T(V)$ si y solo si existen $a,b>0$ tales que
\[
 a\norm v\leq\norm{Tv}\leq b\norm v\qquad(v\in V).
\]
Para concluir que es un isomorfismo de $V$ sobre $W$ hay que exigir además $T(V)=W$.
\end{teorema}
\begin{proof}
Si $T$ es un isomorfismo sobre su imagen y $V\ne\{0\}$, la continuidad da
\[
 \norm{Tv}\leq\norm T\norm v,
 \qquad\norm v=\norm{T^{-1}Tv}\leq\norm{T^{-1}}\norm{Tv}.
\]
Por tanto sirven $a=1/\norm{T^{-1}}$ y $b=\norm T$. Para el espacio nulo pueden elegirse cualesquiera constantes positivas.

Recíprocamente, la cota superior hace continuo a $T$. Si $Tv=0$, la inferior obliga a que $v=0$, por lo que $T$ es inyectivo. Tiene una inversa definida en $T(V)$, y para $w=Tv$,
\[
 \norm{T^{-1}w}=\norm v\leq a^{-1}\norm w.
\]
Esta inversa es continua. No hemos probado que cada $w\in W$ sea de la forma $Tv$, porque tal afirmación no se deduce de las desigualdades.
\end{proof}

\begin{observacion}[Corrección de la lectura original]
La inclusión $T:\R\to\R^2$, $T(t)=(t,0)$, es una isometría y satisface ambas cotas con $a=b=1$, pero no es sobreyectiva. Es un contraejemplo al enunciado original si ``isomorfismo'' se entiende sobre todo el codominio. La formulación anterior distingue la imagen del codominio.
\end{observacion}

La inversa de un isomorfismo es un isomorfismo, y la composición de dos isomorfismos también: la inversa de $ST$ es $T^{-1}S^{-1}$, y la composición de operadores continuos es continua. Cuando $V\ne\{0\}$,
\[
 1=\norm{I_V}=\norm{T^{-1}T}\leq\norm{T^{-1}}\norm T.
\]
El factor $2I$ es un isomorfismo pero no una isometría. Por otra parte, si $V$ es completo y $T$ tiene cota inferior positiva, $T(V)$ es cerrado en $W$: de $Tv_n\to w$ se deduce que $(v_n)$ es de Cauchy; si $v_n\to v$, la continuidad da $w=Tv$.

\section{La extensión de Hahn--Banach}

La construcción de funcionales a partir de una base de Hamel produce funcionales algebraicas, pero no controla su norma. El siguiente argumento conserva una cota durante la extensión. Se presenta con detalle para justificar tanto la identificación isométrica con una parte del bidual como los teoremas de separación de conjuntos convexos que siguen.

\begin{definicion}
En un espacio vectorial real, $p:V\to\R$ es \emph{sublineal} si
\[
 p(x+y)\leq p(x)+p(y),\qquad p(tx)=tp(x)\quad(t\geq0).
\]
No se exige $p(-x)=p(x)$. Una norma es un ejemplo, pero los funcionales asociados a conjuntos convexos pueden no ser simétricos.
\end{definicion}

\begin{teorema}[Hahn--Banach real, forma dominada]\label{ch21:hb-real}
Sea $p$ sublineal en un espacio vectorial real $V$, sea $M\leq V$ y sea $f:M\to\R$ lineal con $f(m)\leq p(m)$ para todo $m\in M$. Existe una funcional lineal $F:V\to\R$ que extiende a $f$ y satisface $F(v)\leq p(v)$ para todo $v\in V$.
\end{teorema}
\index{Hahn--Banach, teorema de}
\begin{proof}
Primero extendemos una sola dimensión. Si $x_0\notin M$, cualquier extensión a $M+\R x_0$ tiene la forma
\[
 f_1(m+tx_0)=f(m)+tc,
\]
donde $c$ es el valor que asignaremos a $x_0$. La representación es única porque $x_0\notin M$. Para que $f_1\leq p$ necesitamos
\[
 \sup_{y\in M}\{f(y)-p(y-x_0)\}\leq c
 \leq\inf_{z\in M}\{p(z+x_0)-f(z)\}.
\]
El intervalo no es vacío. En efecto, para $y,z\in M$,
\[
 f(y)+f(z)=f(y+z)\leq p(y+z)
 \leq p(y-x_0)+p(z+x_0).
\]
Reordenando obtenemos que cada término del extremo izquierdo es menor o igual que cada término del derecho. Además ambos extremos son finitos: tomar $y=0$ proporciona una cota inferior finita para el supremo, y tomar $z=0$ proporciona una cota superior finita; la desigualdad cruzada controla las otras dos cotas. Podemos elegir un número real $c$ en ese intervalo.

Si $t>0$, la cota superior aplicada a $z=m/t$ da
\[
 f(m)+tc\leq t\,p(m/t+x_0)=p(m+tx_0).
\]
Si $t=-s<0$, la cota inferior aplicada a $y=m/s$ da
\[
 f(m)-sc\leq s\,p(m/s-x_0)=p(m-sx_0).
\]
Para $t=0$ se trata de la hipótesis inicial. Hemos construido la extensión de una dimensión.

Considere ahora todos los pares $(N,g)$ donde $M\subseteq N\subseteq V$ es un subespacio, $g$ extiende a $f$ y $g\leq p$ en $N$. Ordénelos por extensión. La familia no es vacía porque contiene $(M,f)$. En una cadena, las funciones coinciden en los dominios comunes. Por ello la unión de sus dominios es un subespacio y la unión de sus funciones es una funcional lineal bien definida dominada por $p$. Es una cota superior de la cadena. Zorn proporciona un par maximal $(N,g)$. Si $N\ne V$, el paso de una dimensión lo extendería a un subespacio mayor, contradiciendo maximalidad. Así $N=V$.
\end{proof}

\begin{corolario}[Extensión con conservación de la norma]\label{ch21:hb-norma}
Si $M$ es un subespacio de un espacio normado real o complejo $V$ y $f\in M^*$, existe $F\in V^*$ con $F|_M=f$ y $\norm F=\norm f$.
\end{corolario}
\begin{proof}
En el caso real aplique el teorema a $p(v)=\norm f\norm v$. Obtenemos $F(v)\leq\norm f\norm v$. Aplicar la misma desigualdad a $-v$ da $-F(v)\leq\norm f\norm v$, luego $|F(v)|\leq\norm f\norm v$ y $\norm F\leq\norm f$. La restricción da la desigualdad contraria.

En el caso complejo considere $V$ y $M$ como espacios reales. La función $g=\operatorname{Re}f$ es real lineal y está dominada por $p(v)=\norm f\norm v$. Extiéndala a una funcional real lineal $G$ dominada por $p$, y defina
\[
 F(v)=G(v)-iG(iv).
\]
Es aditiva y real homogénea. Además,
\[
 F(iv)=G(iv)-iG(-v)=G(iv)+iG(v)=iF(v),
\]
lo que implica homogeneidad compleja. Para $m\in M$, $G(m)=\operatorname{Re}f(m)$ y $G(im)=\operatorname{Re}(if(m))=-\operatorname{Im}f(m)$, así que $F(m)=f(m)$.

Si $F(v)\ne0$, elija un escalar $\lambda$ de módulo $1$ con $\lambda F(v)=|F(v)|$. Entonces
\[
 |F(v)|=\operatorname{Re}F(\lambda v)=G(\lambda v)
 \leq\norm f\norm{\lambda v}=\norm f\norm v.
\]
Si $F(v)=0$, la desigualdad es inmediata. Así $\norm F\leq\norm f$; la restricción vuelve a dar igualdad.
\end{proof}

\begin{corolario}\label{ch21:funcional-normante}
Para cada $v\ne0$ existe $\varphi\in V^*$ con $\norm\varphi=1$ y $\varphi(v)=\norm v$. En consecuencia,
\[\norm v=\sup_{\norm\varphi\leq1}|\varphi(v)|.\]
\end{corolario}
\begin{proof}
En $M=\F v$ defina $f(\alpha v)=\alpha\norm v$. Esta funcional tiene norma $1$, porque $|f(\alpha v)|=\norm{\alpha v}$. Extiéndala por el corolario anterior. La desigualdad $|\varphi(v)|\leq\norm\varphi\norm v$ da una cota superior para el supremo, y la extensión construida alcanza esa cota. Para $v=0$, todos los valores son cero.
\end{proof}

La extensión con conservación de la norma asegura que las funcionales continuas distinguen vectores; la forma dominada de Hahn--Banach permite algo más: adaptar una funcional a la geometría de un conjunto que no tiene por qué ser simétrico. Pasamos ahora de hiperplanos y mediciones a conjuntos convexos, antes de construir sus funcionales de separación.

\section{Conjuntos y combinaciones convexas}\label{lectura:24}

Toda convexidad de estas secciones es \emph{real}. Si el espacio es complejo, trabajamos con su espacio vectorial real subyacente. Los coeficientes de las combinaciones convexas son reales no negativos; las desigualdades de separación se escriben para funcionales reales, o para la parte real de una funcional compleja. No existe un orden de los números complejos que permita escribir sin más $\varphi(a)<\varphi(b)$. La convexidad permite interpolar entre puntos sin salir de un conjunto; una funcional que alcanza allí un máximo permitirá identificar sus partes extremas.



\begin{definicion}
El segmento entre $x$ e $y$ es
\[
 [x,y]=\{(1-t)x+ty:0\leq t\leq1\}.
\]
Un conjunto $A$ es \emph{convexo} si $[x,y]\subseteq A$ para todos $x,y\in A$. Un conjunto es un \emph{cono convexo} si $\alpha x+\beta y$ pertenece a él siempre que $x,y$ pertenezcan a él y $\alpha,\beta\geq0$.
\end{definicion}
\index{convexidad}\index{cono convexo}

\begin{ejemplo}
El vacío y todo subespacio vectorial son convexos. Toda bola abierta de un espacio normado es convexa: si $\norm{y-x}<r$, $\norm{z-x}<r$ y $0\leq t\leq1$, entonces
\[
 \norm{(1-t)y+tz-x}\leq(1-t)\norm{y-x}+t\norm{z-x}<r.
\]
La última desigualdad vale también en $t=0,1$, donde queda una sola de las dos hipótesis. Para la bola cerrada, las mismas estimaciones con $\leq r$ dan $\norm{(1-t)y+tz-x}\leq r$. Un cono convexo es convexo tomando $\alpha=1-t$ y $\beta=t$. Si un cono no es vacío, contiene al cero, porque se permiten ambos coeficientes iguales a cero.
\end{ejemplo}

\begin{proposicion}\label{ch24:interior-cierre}
El interior y el cierre de un conjunto convexo de un espacio normado son convexos.
\end{proposicion}
\begin{proof}
Si $x,y\in\interior A$, elija $r>0$ tal que $B(x,r)$ y $B(y,r)$ estén contenidos en $A$. Para $0\leq t\leq1$ y $\norm h<r$, los puntos $x+h,y+h$ pertenecen a $A$. Por convexidad,
\[
 (1-t)(x+h)+t(y+h)=(1-t)x+ty+h\in A.
\]
Así la bola de radio $r$ alrededor de $(1-t)x+ty$ está en $A$, y ese punto es interior. Si el interior es vacío, su convexidad se cumple por definición.

Para el cierre, si $x,y\in\cl A$, elija sucesiones $x_n,y_n\in A$ convergentes a ellos. Para cada $t\in[0,1]$, los puntos $(1-t)x_n+ty_n$ están en $A$ y convergen a $(1-t)x+ty$ por continuidad de las operaciones. Por ello el límite pertenece a $\cl A$.
\end{proof}

\begin{definicion}
La \emph{envolvente convexa}, o cápsula convexa, de $A$, denotada por $\operatorname{co}(A)$, es la intersección de todos los conjuntos convexos que contienen a $A$. La \emph{envolvente convexa cerrada} es la intersección de todos los conjuntos cerrados convexos que contienen a $A$.
\end{definicion}
\index{envolvente convexa}\index{envolvente convexa!cerrada}

Estas familias no son vacías, porque contienen a todo el espacio. La intersección de conjuntos convexos es convexa: si dos puntos están en todos los conjuntos de la familia, todo su segmento está en cada uno y por tanto en la intersección. La intersección también es cerrada cuando todos los conjuntos lo son. Así las definiciones producen los menores conjuntos con las propiedades indicadas.

\begin{proposicion}\label{ch24:combinaciones}
La envolvente convexa de $A$ es el conjunto de todas sus combinaciones convexas finitas:
\[
 \operatorname{co}(A)=
 \left\{\sum_{j=1}^n\alpha_ja_j:
 n\geq1,\ a_j\in A,\ \alpha_j\geq0,\ \sum_{j=1}^n\alpha_j=1\right\}.
\]
Además, su envolvente convexa cerrada es $\cl{\operatorname{co}(A)}$.
\end{proposicion}
\begin{proof}
Denote por $D$ el conjunto del lado derecho. Contiene a $A$ tomando una sola variable con coeficiente $1$. Si $x=\sum_{j=1}^n\alpha_ja_j$ e $y=\sum_{k=1}^m\beta_kb_k$, entonces
\[
 (1-t)x+ty=\sum_{j=1}^n(1-t)\alpha_ja_j+\sum_{k=1}^m t\beta_kb_k.
\]
Los coeficientes son no negativos y su suma es $(1-t)+t=1$. Por tanto $D$ es convexo, y la minimalidad de $\operatorname{co}(A)$ da $\operatorname{co}(A)\subseteq D$.

Para la otra inclusión, pruebe por inducción que todo conjunto convexo que contiene a $A$ contiene esas combinaciones. Para un término es inmediato. En una combinación de $n+1$ términos, si $\alpha_{n+1}=1$, el resultado es $a_{n+1}$. Si $\alpha_{n+1}<1$, ponga $s=1-\alpha_{n+1}>0$ y
\[
 y=\sum_{j=1}^n\frac{\alpha_j}{s}a_j.
\]
Por la hipótesis inductiva, $y$ pertenece al conjunto convexo, y la combinación original es $sy+\alpha_{n+1}a_{n+1}$, que también pertenece a él. Esto prueba $D\subseteq\operatorname{co}(A)$.

Por la proposición~\ref{ch24:interior-cierre}, $\cl{\operatorname{co}(A)}$ es cerrado, convexo y contiene a $A$. Cualquier cerrado convexo que contenga a $A$ contiene primero a $\operatorname{co}(A)$ y después a su cierre. Así ese cierre es exactamente la envolvente convexa cerrada. Si $A=\varnothing$, todas las envolventes son vacías y las afirmaciones siguen siendo válidas.
\end{proof}

\section{Separación: construcción de la funcional}

Para separar conjuntos convexos, primero convertimos la geometría de un abierto convexo en una función sublineal. La simetría no será necesaria; por eso utilizamos la versión real dominada de Hahn--Banach, no solamente su versión para normas.

\begin{lema}[Funcional de Minkowski]\label{ch24:minkowski}
Sea $U$ abierto convexo en un espacio normado real y suponga $0\in U$. Defina
\[
 p_U(x)=\inf\{t>0:x\in tU\}.
\]
Entonces $p_U$ es finito, no negativo y sublineal. Si $B(0,r)\subseteq U$, entonces $p_U(x)\leq\norm x/r$. Además,
\[U=\{x:p_U(x)<1\}.\]
\end{lema}
\index{Minkowski, funcional de}
\begin{proof}
Como $0$ es interior, existe $r>0$ con $B(0,r)\subseteq U$. Para todo $t>\norm x/r$, el vector $x/t$ está en esa bola y, por tanto, $x\in tU$. Esto prueba finitud y $p_U(x)\leq\norm x/r$. La no negatividad procede de que el ínfimo se toma entre números positivos, y $p_U(0)=0$ porque $0\in tU$ para todo $t>0$.

Para $s>0$, cambiar $t$ por $su$ en la definición da $p_U(sx)=sp_U(x)$; para $s=0$ la igualdad ya está probada. Observe que $0\in U$ y convexidad implican $aU\subseteq bU$ cuando $0<a<b$: si $u\in U$, entonces $(a/b)u+(1-a/b)0\in U$.

Si $a>p_U(x)$, la definición de ínfimo permite elegir $t<a$ con $x\in tU$, y la inclusión anterior da $x\in aU$. Si también $b>p_U(y)$, los puntos $x/a$ e $y/b$ están en $U$. Su combinación convexa
\[
 \frac{x+y}{a+b}=\frac a{a+b}\frac xa+\frac b{a+b}\frac yb
\]
pertenece a $U$. Así $p_U(x+y)\leq a+b$. Haciendo $a\downarrow p_U(x)$ y $b\downarrow p_U(y)$ obtenemos subaditividad.

Si $p_U(x)<1$, existe $0<t<1$ con $x\in tU$, y por convexidad $tU\subseteq U$. Por tanto $x\in U$. Recíprocamente, si $x\in U$ y $x\ne0$, la apertura de $U$ y la continuidad de $s\mapsto sx$ permiten elegir $s>1$ suficientemente próximo a $1$ con $sx\in U$. Entonces $x\in s^{-1}U$, de modo que $p_U(x)\leq s^{-1}<1$. Para $x=0$ ya se sabe que $p_U(0)=0$.
\end{proof}

\begin{observacion}
No hemos afirmado que $p_U(-x)=p_U(x)$. Por ejemplo, para $U=(-1,2)\subseteq\R$, $p_U(1)=1/2$ y $p_U(-1)=1$. La sublinealidad exige homogeneidad solo para escalares no negativos. Esta precisión es esencial para utilizar correctamente la convexidad sin añadir una simetría que el conjunto no tiene.
\end{observacion}

\begin{lema}\label{ch24:separacion-abierto}
Si $D$ es abierto, convexo, no vacío y $0\notin D$, existe una funcional real lineal continua $\ell\ne0$ tal que $\ell(d)>0$ para todo $d\in D$.
\end{lema}
\begin{proof}
Elija $a\in D$ y ponga $U=D-a$. Es abierto convexo y contiene al origen. El vector $x_0=-a$ no pertenece a $U$, porque de lo contrario $0\in D$. Por el lema anterior, $p_U(x_0)\geq1$.

En la recta $\R x_0$ defina $f(tx_0)=t p_U(x_0)$. Para $t\geq0$, la homogeneidad de $p_U$ da $f(tx_0)=p_U(tx_0)$. Para $t<0$, $f(tx_0)<0\leq p_U(tx_0)$. Así $f\leq p_U$ en su dominio. Hahn--Banach real, teorema~\ref{ch21:hb-real}, proporciona una extensión $F$ a todo el espacio con $F\leq p_U$.

Si $B(0,r)\subseteq U$, aplicamos la cota del lema tanto a $x$ como a $-x$ y obtenemos
\[
 F(x)\leq\norm x/r,\qquad -F(x)=F(-x)\leq\norm x/r.
\]
Por ello $F$ es continua. Además $F(x_0)=p_U(x_0)\geq1$, de modo que no es cero. Si $d\in D$, entonces $d-a\in U$ y
\[
 F(d)=F(d-a)+F(a)<1-p_U(x_0)\leq0.
\]
La funcional $\ell=-F$ satisface la conclusión.
\end{proof}

\begin{teorema}[Separación estricta de un compacto y un cerrado]\label{ch24:separacion}
Sean $A,B$ no vacíos, convexos y disjuntos en un espacio normado real, con $A$ compacto y $B$ cerrado. Existe $\ell\in V^*$ tal que
\[
 \sup_{a\in A}\ell(a)<\inf_{b\in B}\ell(b).
\]
En un espacio complejo, existe una funcional compleja continua $\varphi$ para la cual la misma desigualdad vale con $\operatorname{Re}\varphi$ en lugar de $\ell$.
\end{teorema}
\index{separacion@separación!de conjuntos convexos}
\begin{proof}
El conjunto $C=B-A$ es convexo. Es cerrado porque es la suma del cerrado $B$ y del compacto $-A$, como se demostró en el ejercicio~\ref{ej:lectura-19-2}. La disjunción implica $0\notin C$. Por cerradura existe $d=\dist(0,C)>0$. Sea $r=d/2$ y
\[
 D=C+B(0,r).
\]
Es convexo y abierto: la convexidad se comprueba sumando las combinaciones correspondientes, y la apertura procede de que es una unión de traslaciones de la bola abierta. No contiene a cero, porque $c+h=0$ con $\norm h<r$ implicaría $\norm c<d$, imposible.

El lema anterior da $\ell\ne0$ con $\ell(c+h)>0$ para $c\in C$ y $\norm h<r$. La simetría de la bola y la definición de norma de una funcional real implican
\[
 \inf_{\norm h<r}\ell(h)=-r\norm\ell.
\]
Para justificarlo, la cota $\ell(h)\geq-\norm\ell\norm h$ da un sentido; en el otro, elija vectores unitarios donde $|\ell|$ se aproxime a $\norm\ell$, cambie su signo para obtener valores negativos y multiplíquelos por escalares menores que $r$ que tiendan a $r$. De $\ell(c)+\ell(h)>0$ para todos esos $h$ se deduce
\[
 \ell(c)\geq r\norm\ell>0.
\]
En particular, para todo $a\in A$ y $b\in B$,
\[
 \ell(b)-\ell(a)\geq r\norm\ell.
\]
Como $A$ es compacto, $\ell$ alcanza en él un máximo finito. Tomándolo en la desigualdad y después tomando el ínfimo sobre $B$, obtenemos
\[
 \inf_B\ell\geq\max_A\ell+r\norm\ell>\max_A\ell.
\]

En el caso complejo aplique la prueba al espacio real subyacente. A partir de la funcional real $\ell$, defina $\varphi(v)=\ell(v)-i\ell(iv)$. El cálculo del corolario~\ref{ch21:hb-norma} muestra que es compleja lineal y continua, y su parte real es precisamente $\ell$. Esto da la versión compleja sin ordenar números complejos.
\end{proof}

\section{Puntos extremos y el teorema de Krein--Milman}

\begin{definicion}
Un punto $x\in A$ es \emph{extremo} si, cuando $x=(1-t)y+tz$ con $y,z\in A$ y $0<t<1$, necesariamente $y=z=x$. El conjunto de puntos extremos se denota $\operatorname{ext}(A)$.

Un subconjunto $F\subseteq A$ es \emph{extremo en $A$} si toda combinación $(1-t)y+tz\in F$, con $y,z\in A$ y $0<t<1$, obliga a que $y,z\in F$. Si además $F$ es convexo, se denomina una \emph{cara} de $A$.
\end{definicion}
\index{punto extremo}\index{cara!de un convexo}

La definición equivalente de punto extremo dice que no puede escribirse como combinación convexa estricta de dos puntos \emph{distintos} de $A$. Esa palabra falta en la formulación original: sin ella, $x=(1-t)x+tx$ excluiría a todos los puntos. Una singleton $\{x\}$ es extrema en $A$ exactamente cuando $x$ es un punto extremo.

\begin{lema}\label{ch24:caras}
Sea $K$ compacto no vacío y sea $\ell$ una funcional real continua. El conjunto
\[
 F=\{x\in K:\ell(x)=\max_K\ell\}
\]
es no vacío, compacto y extremo en $K$. Si $K$ es convexo, $F$ también lo es. Además, si $G$ es extremo en $F$ y $F$ es extremo en $A$, entonces $G$ es extremo en $A$.
\end{lema}
\begin{proof}
El máximo $m$ existe por compacidad y continuidad. El conjunto $F=K\cap\ell^{-1}(\{m\})$ es un cerrado no vacío de $K$, luego compacto. Si $(1-t)x+ty\in F$ con $x,y\in K$ y $0<t<1$, entonces
\[
 0=m-\ell((1-t)x+ty)
 =(1-t)(m-\ell(x))+t(m-\ell(y)).
\]
Los dos sumandos son no negativos y sus coeficientes son estrictamente positivos. Por ello ambos déficits son cero, de modo que $x,y\in F$. Si $K$ es convexo y $x,y\in F$, el segmento pertenece a $K$ y la linealidad da $\ell((1-t)x+ty)=m$, por lo que está en $F$.

Para la transitividad, si $(1-t)x+ty\in G\subseteq F$ y $x,y\in A$, la extremalidad de $F$ fuerza $x,y\in F$; después la extremalidad de $G$ dentro de $F$ fuerza $x,y\in G$.
\end{proof}

\begin{teorema}[Krein--Milman para compactos en norma]\label{ch24:krein-milman}
Si $A$ es no vacío, compacto y convexo en un espacio normado, entonces $\operatorname{ext}(A)\ne\varnothing$ y
\[
 A=\cl{\operatorname{co}(\operatorname{ext}(A))}.
\]
\end{teorema}
\index{Krein--Milman, teorema de}
\begin{proof}
Primero encontraremos un punto extremo. Sea $\mathcal F$ la familia de los subconjuntos \emph{no vacíos}, compactos y extremos de $A$. No es vacía porque $A\in\mathcal F$. Ordénela por inclusión inversa: $F\preceq G$ significa $G\subseteq F$. Buscamos así un elemento maximal para ese orden, que será un conjunto mínimo por inclusión.

Sea $\mathcal C$ una cadena no vacía de $\mathcal F$. Sus miembros son compactos y, por tanto, cerrados en el compacto $A$. Cualquier intersección finita de miembros de una cadena coincide con el más pequeño de esos miembros y no es vacía. La propiedad de intersección finita de los compactos implica
\[
 K=\bigcap_{F\in\mathcal C}F\ne\varnothing.
\]
Es compacto por ser cerrado en $A$. Es extremo porque, si una combinación estricta de puntos de $A$ pertenece a $K$, pertenece a cada $F\in\mathcal C$; la extremalidad de cada $F$ obliga a que ambos extremos pertenezcan a todos ellos y, por tanto, a $K$. Así $K\in\mathcal F$ y es una cota superior de la cadena en el orden inverso. Para la cadena vacía, cualquier miembro de $\mathcal F$ es una cota superior. Zorn proporciona un miembro minimal por inclusión $F_0$.

Si $F_0$ contuviera dos puntos distintos $u,v$, Hahn--Banach real aplicado al vector $v-u$ proporcionaría una funcional real continua $\ell$ con $\ell(v-u)=\norm{v-u}>0$. Por tanto $\ell(u)<\ell(v)$. El subconjunto de máximos de $\ell$ en $F_0$ es compacto, no vacío y extremo en $F_0$ por el lema anterior. Por transitividad, es extremo en $A$. Es un subconjunto propio de $F_0$ porque no contiene a $u$, contradiciendo la minimalidad. Así $F_0=\{e\}$ y $e\in\operatorname{ext}(A)$.

Ponga $E=\operatorname{ext}(A)$ y $C=\cl{\operatorname{co}(E)}$. Como $A$ es cerrado convexo y contiene a $E$, se tiene $C\subseteq A$. En particular, $C$ es compacto por ser cerrado en $A$; no estamos usando una afirmación general de compacidad de envolventes convexas.

Si existiera $x\in A\setminus C$, el teorema de separación aplicado a $C$ y $\{x\}$ daría una funcional real continua $\ell$ con
\[
 \sup_C\ell<\ell(x)\leq\max_A\ell=:m.
\]
La cara $F=\{a\in A:\ell(a)=m\}$ es no vacía, compacta y convexa, y es extrema en $A$. La primera parte de la demostración, aplicada a $F$, proporciona un punto extremo $e$ de $F$. La transitividad del lema hace que $e$ sea extremo en $A$, así que $e\in E\subseteq C$. Pero $\ell(e)=m>\sup_C\ell$, contradicción. En consecuencia $A\subseteq C$, y queda probada la igualdad.
\end{proof}

El cierre en el enunciado permite límites de combinaciones convexas finitas; no equivale a permitir una suma infinita arbitraria sin comprobar su convergencia. La versión aquí demostrada usa compacidad para la norma. Otras topologías conducen a versiones más generales, pero requieren hipótesis adicionales que no se están utilizando en este capítulo.

\section{Ejercicios propuestos: dualidad, distancia e invertibilidad}

\begin{ejercicio}[Lectura 21, ejercicio 1]\label{ej:lectura-21-1}
Demuestre que el espacio de todas las funcionales lineales de $\F^n$ es isomorfo a $\F^n$.
\end{ejercicio}


\begin{ejercicio}[Lectura 21, ejercicio 2]\label{ej:lectura-21-2}
Sea $\varphi\in V^*$ y $M=\ker\varphi$. Pruebe que
\[|\varphi(v)|=\norm\varphi\,\dist(v,M)\qquad(v\in V).\]
\end{ejercicio}


\begin{ejercicio}[Lectura 21, ejercicio 3]\label{ej:lectura-21-3}
Sea $A\subseteq V$ no vacío y $U=\overline B(0,1)$. Pruebe
\[
 \cl{A+U}=\{v\in V:\dist(v,A)\leq1\}.
\]
Deduzca que $A+U$ es cerrado si y solo si todo $v$ con $\dist(v,A)=1$ admite $a\in A$ con $\norm{v-a}=1$.
\end{ejercicio}


\begin{ejercicio}[Lectura 21, ejercicio 4]\label{ej:lectura-21-4}
Sea $\norm\varphi=1$, $M=\ker\varphi$ y $U=\overline B(0,1)$. Demuestre
\[
 \cl{M+U}=\{v:|\varphi(v)|\leq1\}.
\]
Pruebe que $M+U$ es cerrado si y solo si existe $u\in U$ con $|\varphi(u)|=1$.
\end{ejercicio}


\begin{ejercicio}[Lectura 21, ejercicio 5]\label{ej:lectura-21-5}
Sea $V$ de Banach y $T\in\Bcal(V,V)$ con $\sum_{n=0}^{\infty}\norm{T^n}<\infty$. Demuestre que $I-T$ es un isomorfismo y calcule su inversa.
\end{ejercicio}

\index{Neumann, serie de}

\begin{ejercicio}[Lectura 21, ejercicio 6]\label{ej:lectura-21-6}
Sean $V,W$ de Banach, con $V\ne\{0\}$, y $T:V\to W$ un isomorfismo. Si $S\in\Bcal(V,W)$ satisface $\norm{S-T}<1/\norm{T^{-1}}$, demuestre que $S$ es un isomorfismo. Deduzca que los isomorfismos forman un abierto de $\Bcal(V,W)$, incluyendo por separado el caso $V=W=\{0\}$.
\end{ejercicio}


\begin{ejercicio}[Lectura 21, ejercicio 7]\label{ej:lectura-21-7}
Sea $T$ un isomorfismo de $V$ sobre un subespacio de $W$. Pruebe que, si $\norm{T'-T}$ es suficientemente pequeña, $T'$ es un isomorfismo sobre su imagen.
\end{ejercicio}


\begin{ejercicio}[Lectura 21, ejercicio 8]\label{ej:lectura-21-8}
Pruebe que la evaluación canónica
\[
 J:V\longrightarrow V^{**},\qquad J(v)(\varphi)=\varphi(v),
\]
es una isometría lineal.
\end{ejercicio}


\section{Ejercicios propuestos: envolventes y separación}

\begin{ejercicio}[Lectura 24, ejercicio 1]\label{ej:lectura-24-1}
Demuestre que la intersección de una familia no vacía de conjuntos convexos es convexa.
\end{ejercicio}


\begin{ejercicio}[Lectura 24, ejercicio 2]\label{ej:lectura-24-2}
Si $A,B\subseteq V$ y $C\subseteq W$ son convexos, $\alpha\in\R$ y $T:V\to W$ es lineal, pruebe que $\alpha A$, $A+B$, $T(A)$ y $T^{-1}(C)$ son convexos.
\end{ejercicio}


\begin{ejercicio}[Lectura 24, ejercicio 3]\label{ej:lectura-24-3}
Sean $A_1,\ldots,A_n$ conjuntos convexos no vacíos. Demuestre que $\operatorname{co}(A_1\cup\cdots\cup A_n)$ está formado por las combinaciones $\sum_{j=1}^n\alpha_jx_j$, con $x_j\in A_j$, $\alpha_j\geq0$ y $\sum_j\alpha_j=1$.
\end{ejercicio}


\begin{ejercicio}[Lectura 24, ejercicio 4: revisión y solución]\label{ej:lectura-24-4}
La lectura propone, para conjuntos cerrados $A,B$, las afirmaciones: (i) si ambos son compactos, $\operatorname{co}(A\cup B)$ es compacto; (ii) si $A$ es compacto y $B$ acotado, $\operatorname{co}(A\cup B)$ es cerrado. Examine ambas afirmaciones y pruebe sus versiones correctas. Demuestre además que la envolvente convexa de $\{(0,0)\}\cup\{(1,y):y\in\R\}$ no es cerrada.
\end{ejercicio}


\begin{ejercicio}[Lectura 24, ejercicio 5]\label{ej:lectura-24-5}
Si $A$ es convexo, demuestre que $C=\bigcup_{s\geq0}sA$ es un cono convexo.
\end{ejercicio}


\begin{ejercicio}[Lectura 24, ejercicio 6]\label{ej:lectura-24-6}
Demuestre que la envolvente convexa cerrada de un compacto no vacío $A$ es el cierre de su envolvente convexa. Determine si la compacidad es necesaria.
\end{ejercicio}


\begin{ejercicio}[Lectura 24, ejercicio 7]\label{ej:lectura-24-7}
Sean $A,B$ convexos y disjuntos, y sea $x\in V$. Pruebe que
\[
 \operatorname{co}(A\cup\{x\})\cap B=\varnothing
 \quad\text{o}\quad
 \operatorname{co}(B\cup\{x\})\cap A=\varnothing.
\]
\end{ejercicio}
